\section{Roots of Western Civilisation}

Among the many would-be ``saviours'' of Western civilisation, there are precious few who know exactly what they are supposed to preserve. For example, many prefer the values of the Enlightenment and praise scientific discovery; while that may be, in a restricted sense, one of the flowers, it is hardly the root.

To locate the roots, the first source would be the works of George Dumezil, who identified the common social and religious structures of Indo-European peoples from Vedic India to Ireland, and all points in-between. Even if you deny that this was the work of a particular race, it cannot be denied that it was the manifestation of a particular spirit.

In \emph{Sintesi di dottrina dell razza}, Julius Evola mentioned the qualities ``characteristic of the great Aryan civilisations of the Orient, ancient Rome, up to the Roman-Germanic Middle Ages''. So to understand the roots of Western civilisation, it is necessary to return to those sources: the Vedic civilisation of India, the classical civilisations of Greece and Rome, and finally the Catholic Middle Ages as created by the Germanic and classical Roman currents.

Let us not forget, there was no ``Europe'' until the Middle Ages. Ancient Rome was built around the Mediterranean and included what is now called the Near East and the Mahgreb. Savitri Devi in \emph{Son of the Sun} writes:

\begin{quotex}
If we consider the Western world as a whole (Europe \emph{and its background}), and not only the small portion of it which one generally has in mind when speaking of ``the West'', then we have to include in it the countries of the Bible---Syria, Egypt, Arabia, Iraq---no less than Greece; for they are the geographical and cultural background of Christianity, the religion of Europe for centuries.
\end{quotex}


After the Middle Ages, Western civilisation takes a different turn and, no longer nourished by its roots, its decline commenced.


\flrightit{Posted on 2009-10-13 by Cologero}

\begin{center}* * *\end{center}

\begin{footnotesize}
\begin{sffamily}
\texttt{kadambari on 2010-08-04 at 22:31 said:}

Sorry to write so many comments (I have really enjoyed reading some of the selections on this site!!!)-- this post was very nice until it mentions Savitri Devi who was a nut job!

I can't imagine anyone with intelligence can even take this charlatan of a woman seriously! What do you think of this piece by Elst on her?\\
\url{http://koenraadelst.bharatvani.org/articles/SavitriBouchet.html}


\hfill

\texttt{Cologero on 2010-08-04 at 23:15 said:}

We are not interested in gossip or personalities at Gornahoor. Many Gornahoor readers also read Devi, so I quoted her to make a point, which is that the so-called ``West'' includes more than Northern Europe. Please restrict your comments to the particular passage. Even if you disagree with it, it doesn't seem particularly irrational to me.

\hfill

\texttt{kadambari on 2010-08-05 at 08:41 said:}

Certainly, I guess, by virtue of adopting the Mid-East religions, Europe is certainly closer to those areas that those of use farther East who do not have anything in common with them whatsoever! A small educated fast disappearing Parsi community retains their original Aryan religion fleeing conversion and maintains the ancient fire alive... . However their outlook is compatible with the native religions in our areas and I have many Parsi friends. As for Savitri Devi, she should just have stuck perhaps to her own mythologies and not dabble in those of other people's which she clearly does not understand... \hfill

\texttt{kadambari on 2010-08-05 at 09:20 said:}

But Europe had much in common with those areas when Syria and Lebanon were Greek colonies and were thriving intellectual centers, and even later when they became centers of the Christian faith; they are Mohammedian now. The so called Golden Age of Islam was mostly Persians contributing their native genius and writing in Arabic after their conversion... A great deal of the knowledge of the Greeks and Hindus was also absorbed and copied by the Arabs. Also much of central Asia was Zoroasrtian, Buddhist as well as Hindu in Afghanistan. Do you think those people even have a memory of their non-Islamic past? I do not think so. After the seventh centuries those areas became different in outlook as a consequence of religion. Does Europe really have much in common with these areas today? I think it might have more in common with a developed country like Japan although they are Buddhist. In the Middle East, I hear most of the Christian communities in Lebanon and the Holy Land are fleeing due to the violence... .Just some thoughts... \hfill

\end{sffamily}
\end{footnotesize}

